\section{Conclusions and further research questions}
\label{sec:conclusion}

\subsection{Main conclusion}

The broad mathematical content of the repository does not require a separate foundational method for every subject.  A large common core is visible:

\begin{enumerate}
\item specialize a fixed construction to primitive-recursive syntax;
\item propagate a terminal request backward to a finite demand;
\item execute exact loops, bounded searches, and certificate checks forward;
\item keep non-unique data relational;
\item eliminate choices by finite viability or stable observation;
\item demand quantitative guards at discontinuities and partial operations;
\item encode the whole finite execution by a quantifier-free invariant;
\item return an actual \(\Prf\ \PRAAx\) derivation.
\end{enumerate}

The PFIG normal form turns this into one reusable backend theorem.  Its most important effect is methodological: it forces every hidden appeal to choice, compactness, exact zero testing, a convergence modulus, or a universal evaluator to appear as an explicit interface obligation.

\subsection{What has been resolved}

At the human mathematical level, this report gives complete arguments for:

\begin{itemize}
\item structural Booleanization of fixed bounded formulas;
\item least bounded search and finite choice;
\item proof-producing deterministic traces;
\item reverse demand and fixed-graph compilation;
\item finite relational viability and least viable paths;
\item coherent regular names from shrinking viable clouds;
\item exact and metric stable-observer elimination;
\item branch stabilization under a gap;
\item finite-cover gluing and representative-free global codes;
\item equivalence of explicit height bounds and complete finite lists;
\item negative translation and retained classical-to-intuitionistic transfer.
\end{itemize}

The report also explains why several unrestricted requests should be reformulated: exact root labels at collisions, exact axes at the identity, a universal PR evaluator, a Cauchy modulus from metastability alone, exact branch data from arbitrary approximate coefficients, and search cutoffs from mere finiteness.

\subsection{Immediate formalization questions}

\begin{enumerate}
\item What is the smallest object-iPRA Boolean library sufficient to reflect all bounded formulas used by current packets?  A syntax-directed compiler may be less work than a large hand-written arithmetic library.
\item Can the existing \texttt{ObjectIteration}, \texttt{ObjectComposition}, and \texttt{ObjectObservation} code be moved into a stable shared namespace without duplicating the upstream kernel?
\item Which state/table encoding minimizes de Bruijn substitution complexity: rectangular arrays, nested pairs, or bounded-sequence codes?
\item Should the first graph compiler emit specialized Lean files, or should a typed PFIG AST be internalized immediately?  The generated-file route appears lower risk.
\item Can proof terms for Boolean table reflection be compressed by using the native scalar calculus earlier, while preserving a readable ordinary-iPRA API?
\item How much of the negative-translation proof can be automated by generic natural-deduction lemmas independent of arithmetic?
\end{enumerate}

\subsection{Mathematical research questions}

\begin{enumerate}
\item Which classical theorems in the repository have retained conclusions after their existential content is converted to explicit PFIG witnesses?
\item Can root and spectral packets be uniformly reorganized around cluster relations and stable symmetric observables, with exact labels as a guarded corollary?
\item Which compactness arguments admit primitive-recursive finite-net producers with sharp enough moduli for practical formalization?
\item For which Diophantine families can explicit height producers be extracted in iPRA, thereby converting abstract finiteness packets into complete enumerators?
\item Can finite carrier relations and contiguity provide a uniform route from computational topology packets to object-iPRA homology theorems?
\item Is there a useful hierarchy of PFIG interfaces by logical strength---exact, bounded relational, metastable, regular-name, and guarded exact---that predicts which theorem transformations are possible?
\item Can proof mining systematically extract demand maps, tail budgets, and finite certificate formats from existing classical proofs?
\item Which error-budget schedules yield the smallest primitive-recursive state and proof terms for analysis applications?
\end{enumerate}

\subsection{Engineering research questions}

\begin{enumerate}
\item Can a small declarative contract language generate both Python reference code and Lean object formulas from one schema?
\item Can the compiler emit a theorem ledger automatically, marking every imported local certificate and every remaining producer assumption?
\item Can proof-carrying traces be shared across packets so that a checked determinant, polynomial, interval, or finite-net library becomes reusable?
\item Can native-calculus proof checking be used as a compact artifact format after ordinary-iPRA proofs are generated?
\item What is the best strategy for incremental proof checking when a graph contains large finite certificates: direct derivation trees, reflected Boolean computation, or a verified certificate VM specialized to a fixed instruction set?
\end{enumerate}

\subsection{Recommended repository consolidation}

A practical next repository step is to create a shared directory.
\begin{center}
\path{Logic/PrimitiveRecursiveArithmetic/PFIG/}
\end{center}
The directory should contain the modules from \cref{sec:lean-plan}.  Existing framework packets should then be treated as design documents and examples rather than copied kernels.  The theorem-packet README template should require:

\begin{itemize}
\item pinned upstream commit;
\item exact object formulas;
\item status of every \(\Prf\ \PRAAx\) target;
\item test commands and logs;
\item quantitative producer assumptions;
\item explicit barrier statements for intentionally unavailable selectors.
\end{itemize}

This would transform a collection of parallel proposals into one cumulative formalization program.

\subsection{Final perspective}

The right question is not whether ordinary mathematics can be copied verbatim into iPRA.  It is which finite observations, certificates, moduli, and stable relations constitute the proof's actual computational content.  Once those interfaces are made explicit, iPRA is strong enough to verify a surprisingly broad range of mathematics by primitive recursion and quantifier-free induction.  Where the interfaces cannot be supplied, the failure is informative: it identifies a discontinuity, a missing quantitative theorem, or a genuinely non-effective principle rather than an unexplained formalization impasse.
